\section{Experimental Setup}
\label{sec:setup}

\paragraph{Model.} We use XLM-RoBERTa-base \citep{conneau2020xlmr} as the shared encoder for both query and document sides, followed by a single $768 \times 128$ linear projection and per-token $L_2$ normalization. Input sequences are truncated to $64$ tokens. This is a faithful adaptation of the ColBERT architecture to a multilingual backbone, identical for vanilla and IsoColBERT.

\paragraph{Training data.} We train on the English-Spanish parallel split of OPUS-100 \citep{zhang2020opus100}, streamed from the source. EN-ES is the only language pair used for contrastive training; the six other evaluation pairs are zero-shot.

\paragraph{Optimization.} We train for $2000$ optimizer steps at an effective batch size of $32$ (micro-batch $8$, gradient accumulation $4$) with AdamW \citep{loshchilov2019adamw}, learning rate $2 \times 10^{-5}$, weight decay $0.01$, and a linear warmup of $200$ steps. The InfoNCE temperature is $\tau = 0.07$, matching standard contrastive setups. We clip gradient norms at $1.0$.

\paragraph{Evaluation.} We evaluate on FLORES+, the multilingual evaluation continuation of FLORES-200 \citep{nllb2024flores}. FLORES+ consists of parallel sentence translations across 200 languages, so our setup measures \emph{cross-lingual sentence retrieval} rather than retrieval over a heterogeneous document corpus; the candidate pool size is also small relative to operational multilingual retrieval benchmarks such as MIRACL \citep{zhang2023miracl} or mMARCO \citep{bonifacio2021mmarco}. This is a controlled regime in which the geometric effect we study can be measured cleanly, and we discuss the scope implications in Section~\ref{sec:discussion}. For each target language, we combine the FLORES+ \texttt{dev} and \texttt{devtest} splits into a single retrieval pool of approximately $2009$ candidates. We evaluate on seven language pairs: EN-ES as the in-training pair, EN-FR and EN-DE as high-resource zero-shot pairs, and EN-SW, EN-AR, EN-NE (Nepali), and EN-TA (Tamil) as low-resource zero-shot pairs. The latter four all sit far below Spanish, French, and German in XLM-R's pretraining mix; Swahili and Nepali in particular have only a few gigabytes of CommonCrawl text compared to tens of gigabytes for the high-resource pairs \citep{conneau2020xlmr}, and represent the regime where pretrained signal is scarce.

\paragraph{Metrics.} The primary retrieval metric is MaxSim Precision@1 against the full per-language retrieval pool: we score every (query, candidate) pair using \eqref{eq:maxsim} and check whether the true match is rank $1$. For geometry, we measure token intra-cosine \eqref{eq:intracos} over $2048$ randomly drawn valid token vectors per language and effective rank \eqref{eq:erank} over $4096$ stacked token vectors per language. Following \citet{ethayarajh2019contextual} we use the projected, $L_2$-normalized token vectors output by the encoder for both diagnostics.

\paragraph{Seeds and effect-size markers.} All retrieval and geometry numbers are reported as mean $\pm$ standard deviation across three independent random seeds $\{42, 1337, 2024\}$, each producing an independently trained model. We mark deltas with effect-size cutoffs (\textbf{***}~$\Delta\mu > 2\sigma$, \textbf{**}~$\Delta\mu > \sigma$, \textbf{*}~$\Delta\mu > \sigma/2$), where $\sigma$ is the larger of the two seed-wise standard deviations. These are not p-values from a formal statistical test: with three seeds, a paired test has only two degrees of freedom, too few to support conventional significance thresholds. We use the markers as a compact summary of how a mean delta compares to its seed-wise variance, and we additionally list per-seed differences in Section~\ref{sec:results-retrieval} so that sign-consistency across seeds can be checked directly. The hyperparameter sweep over $\lambda$ in Section~\ref{sec:ablations-lambda} reports three-seed means for $\lambda \in \{0, 0.1, 0.5\}$; $\lambda{=}0.5$ is selected as the headline configuration.

\paragraph{Compute.} A single training run takes approximately one GPU-hour on an NVIDIA A100, leading to roughly six GPU-hours for the main multi-seed experiment (three seeds, vanilla and IsoColBERT) and a comparable budget for the ablations.

\paragraph{Code and reproducibility.} All training and evaluation notebooks (one per experiment in Sections~\ref{sec:results} and~\ref{sec:ablations}), the figure-generation scripts, the LaTeX source of this report, and the built PDF are available at \url{https://github.com/eyesacchang/grad-nlp-project}.
